\documentclass[11pt,a4paper]{article}

\usepackage[T1]{fontenc}
\usepackage[utf8]{inputenc}
\usepackage[turkish]{babel}
\usepackage{geometry}
\usepackage{booktabs}
\usepackage{tabularx}
\usepackage{array}
\usepackage{longtable}
\usepackage{xcolor}
\usepackage{fancyhdr}
\usepackage{titlesec}
\usepackage{hyperref}
\usepackage{listings}
\usepackage{enumitem}
\usepackage{parskip}
\usepackage{tcolorbox}

\geometry{
	a4paper,
	top=2.5cm,
	bottom=2.5cm,
	left=2.5cm,
	right=2.5cm,
	headheight=15pt
}

\definecolor{codeblue}{RGB}{32,78,145}
\definecolor{codegray}{RGB}{245,246,248}
\definecolor{codegreen}{RGB}{35,120,80}
\definecolor{codered}{RGB}{170,45,45}
\definecolor{headerbg}{RGB}{238,243,250}

\titleformat{\section}{\large\bfseries\color{codeblue}}{\thesection}{1em}{}[\titlerule]
\titleformat{\subsection}{\normalsize\bfseries}{\thesubsection}{1em}{}

\pagestyle{fancy}
\fancyhf{}
\lhead{\small CodeAlchemist -- 9. Hafta Teknik Dokuman}
\rhead{\small 9 Nisan 2026}
\cfoot{\small CodeAlchemist Team \quad Sayfa \thepage}

\lstset{
	backgroundcolor=\color{codegray},
	basicstyle=\ttfamily\footnotesize,
	keywordstyle=\color{codeblue}\bfseries,
	commentstyle=\color{codegreen},
	stringstyle=\color{codered},
	numbers=left,
	numberstyle=\tiny\color{gray},
	frame=single,
	framerule=0.5pt,
	breaklines=true,
	showstringspaces=false,
	xleftmargin=1.2em,
	xrightmargin=0.3em,
	tabsize=2
}

\tcbuselibrary{skins,breakable}
\newtcolorbox{notebox}[1][] {
	colback=headerbg,
	colframe=codeblue!60,
	fonttitle=\bfseries\small,
	title={#1},
	breakable
}

\hypersetup{
	colorlinks=true,
	linkcolor=codeblue,
	urlcolor=codeblue,
	pdftitle={CodeAlchemist 9. Hafta Teknik Dokuman},
	pdfauthor={CodeAlchemist Team}
}

\begin{document}

\begin{titlepage}
	\centering
	\vspace*{3cm}
	{\Huge\bfseries\color{codeblue} CodeAlchemist\par}
	\vspace{0.5cm}
	{\Large 9. Hafta Teknik Dokuman\par}
	\vspace{0.3cm}
	{\large VS Code Extension Yol Haritasi ve Token Satin Alma Mekanizmasi\par}
	\vspace{1.8cm}
	\hrule height 1pt
	\vspace{0.4cm}
	{\normalsize Uygulama Adimlari \quad $\bullet$ \quad API Sozlesmeleri \quad $\bullet$ \quad Odeme Akisi\par}
	\vspace{0.4cm}
	\hrule height 1pt
	\vfill
	{\normalsize\textbf{Tarih:} 9 Nisan 2026\par}
	\vspace{0.2cm}
	{\normalsize\textbf{Versiyon:} 9.1\par}
	\vspace{0.2cm}
	{\normalsize\textbf{Kapsam:} VS Code extension MVP + Guvenlik ve dayaniklilik iyilestirmeleri\par}
\end{titlepage}

\section{VS Code Extension Mevcut Durum Ozeti}

\begin{notebox}[Kod Tabanina Gore Simdiki Durum]
Rapor yazim anindaki VS Code extension, tek bir is akisi icinde su bilesenleri birlestirmektedir: SecretStorage tabanli API key yonetimi, model ve agent mode secimi, workspace baglam toplama, SSE stream okuma, diff onizleme, action uygulama ve health izleme.
\end{notebox}

\subsection{Bugunku Calisma Haritasi}

\begin{table}[h]
\centering
\begin{tabularx}{\textwidth}{@{}l l X@{}}
\toprule
Katman & Durum & Kod Tabanindaki Karsiligi \\
\midrule
Manifest ve komut yuzeyi & Aktif & \texttt{package.json} icinde \texttt{codeAlchemist.setApiKey}, \texttt{codeAlchemist.ask}, \texttt{codeAlchemist.selectModel}, \texttt{codeAlchemist.toggleAgentMode}, \texttt{codeAlchemist.applyChanges}, \texttt{codeAlchemist.focusChat} komutlari tanimli. \\
Kimlik ve ayar yonetimi & Aktif & API anahtari \texttt{SecretStorage} uzerinden saklaniyor; eski settings degerleri ilk kullanima migre ediliyor; endpoint ve model ayarlari merkezden okunuyor. \\
Workspace baglam toplama & Aktif & Secili kod, acik dosyalar, aktif dosya, calisma alani snapshot'i ve dosya sinirlari tek payload icinde toplanip backend'e gonderiliyor. \\
Istek orkestrasyonu & Aktif & \texttt{/v1/ask} icin zenginlestirilmis JSON govde hazirlaniyor; model, agent modu, dosya yolu ve istemci yetenekleri payload'a ekleniyor. \\
Stream ve yanit akisi & Aktif & SSE parcaciklari Output paneline anlik akiyor; JSON fallback korunuyor; \texttt{[DONE]} sinyali ve hata durumlari guvenli sekilde ele aliniyor. \\
Action uygulama ve guvenlik & Aktif & \texttt{edit\_file} ve \texttt{multi\_edit} aksiyonlari diff onizlemesi ile onaylanarak uygulanir; path safety ve trusted file kontrolleri devrede. \\
Operasyonel izleme & Aktif & Status bar, health monitor, sidebar chat provider ve agent trace ciktilari ayni oturum icinde calisiyor. \\
\bottomrule
\end{tabularx}
\caption{VS Code extension icin mevcut durum ozeti}
\end{table}

\subsection{Ilerleme Asamalari}

Extension tarafindaki ilerleme, sadece tek bir ozellik eklenmesi olarak degil, birbirini tamamlayan katmanlar halinde ilerlemistir. Kod tabanindaki her asama, bir oncekinin uzerine eklenerek daha guvenli ve daha izlenebilir bir is akisi olusturmustur.

\begin{enumerate}[leftmargin=1.5em]
\item \textbf{Asama 1 -- Manifest ve komut yuzeyinin kurulmasi}
\begin{itemize}[leftmargin=1.5em]
\item \texttt{package.json} tarafinda extension'in yuze acilan komutlari, ayar anahtarlari ve activation event'leri tanimlandi.
\item Chat sidebar, API key ayari, model secimi ve agent modu toggle gibi kullanici islemleri komut bazli hale getirildi.
\item Bu asama, extension'in sadece kod calistiran bir parcadan cikip VS Code icinde dogrudan kullanilabilen bir yuzeye donusmesini sagladi.
\end{itemize}

\item \textbf{Asama 2 -- Guvenli durum ve kimlik yonetimi}
\begin{itemize}[leftmargin=1.5em]
\item API key saklama noktasi SecretStorage'a tasindi ve legacy settings kayitlari ilk kullanicida otomatik migrate edilecek sekilde korunur hale getirildi.
\item Endpoint, model ve agentMode gibi runtime ayarlari tek kaynakli hale getirildi; user state ile config state birbirinden ayrildi.
\item Bu asama, anahtar sizintisi riskini azalttigi gibi, extension'i yeniden baslatmadan calisan bir ayar modeline yaklastirdi.
\end{itemize}

\item \textbf{Asama 3 -- Workspace baglaminin zenginlestirilmesi}
\begin{itemize}[leftmargin=1.5em]
\item Aktif editor, secili kod, acik sekmeler ve filtrelenmis workspace snapshot'i birlikte toplanarak backend'e gonderilecek hale getirildi.
\item Dosya turu, yol guvenligi ve trust id uretimi gibi alanlar sayesinde backend tarafinda hangi verinin nereden geldigi izlenebilir oldu.
\item Bu katman, modelin sadece soru degil, proje baglami ile birlikte cevap vermesini saglayan temel altyapiyi olusturdu.
\end{itemize}

\item \textbf{Asama 4 -- Ask akisinin orkestrasyonu}
\begin{itemize}[leftmargin=1.5em]
\item Kullanici sorusu, secili kod ve workspace metadata tek bir zenginlestirilmis \texttt{AskRequestPayload} icinde birlestirildi.
\item \texttt{codeAlchemist.ask} akisi, backend baglantisini dogrulama, payload hazirlama, response okuma ve Output paneline yazma sorumluluklarini net sekilde ayirdi.
\item Bu asama, extension'i bir chat kutusu olmaktan cikartip baglam farkindaligi olan bir islem orkestratorune donusturdu.
\end{itemize}

\item \textbf{Asama 5 -- SSE stream, action parse ve diff onizleme}
\begin{itemize}[leftmargin=1.5em]
\item SSE parcaciklari anlik olarak Output kanalina yazdirilacak sekilde akitildi; bu sayede kullanici cevabi son satira kadar beklemek zorunda kalmadi.
\item Modelin urettigi \texttt{edit\_file} ve \texttt{multi\_edit} aksiyonlari parse edilerek diff onizleme akisi ile birlikte ele alindi.
\item Path safety, trusted path kontrolu ve kullanici onayi, otomatik yazma islerini guvenli sinirlarda tuttu.
\end{itemize}

\item \textbf{Asama 6 -- Sidebar, health monitor ve operasyonal gozlem}
\begin{itemize}[leftmargin=1.5em]
\item Chat sidebar, model secimi, ask gonderimi ve trace goruntuleme icin merkezi kontrol noktasi haline getirildi.
\item Health monitor, backend durumunu izleyerek kullaniciya anlik baglanti sinyali verir hale geldi.
\item Bu son katman, extension'in artik sadece istek atan degil, kendi sagligini ve calisma fazini takip eden bir araca donustugunu gosteriyor.
\end{itemize}
\end{enumerate}

\subsection{Simdiki Urun Pozisyonu}

Kod tabani bugun su noktada duruyor:
\begin{itemize}[leftmargin=1.5em]
\item Kullanici API key, model ve agent modu ayarlarini tek yerden yonetebiliyor.
\item Secili kod ve workspace snapshot'i ile baglamli soru-cevap akisi calisiyor.
\item SSE destekli yanit akisi ve action uygulama onizlemesi aktif.
\item Agent trace, health monitor ve sidebar kontrol yuzeyi birbiriyle entegre halde.
\end{itemize}

Bu nedenle extension, rapor yazim aninda ilk taslak seviyesinde degil; dogrudan kullanilabilen, akisi ayrik ama birbirine bagli modullerden olusan calisir bir MVP durumundadir.

\section*{Yonetici Ozeti}
\addcontentsline{toc}{section}{Yonetici Ozeti}

Bu rapor, iki kritik eksenin teknik durumunu birlestirir:
\begin{itemize}[leftmargin=1.5em]
	\item VS Code extension tarafinda kullanicinin secili kod ile soru gonderebildigi, API anahtari ile dogrulanan ve ciktiyi Output panelinde gosteren MVP akisinin tamamlanmasi.
	\item Uygulama ici token satin alma tarafinda paket listeleme, Stripe checkout olusturma, webhook ile bakiye guncelleme ve islem gecmisi gorunurlugunun devreye alinmasi.
\end{itemize}

Bu iki eksen birlikte ele alindiginda, urun hem editor-ici kullanim (extension) hem de ekonomik surdurulebilirlik (token ekonomisi) acisindan calisir bir omurga kazanmistir.

\newpage
\tableofcontents
\newpage

\section{Kapsam ve Referans Noktalari}

Hafta 9 kapsaminda raporlanan ana dosyalar:
\begin{itemize}[leftmargin=1.5em]
	\item \texttt{vscode-extension/src/extension.ts}
	\item \texttt{vscode-extension/package.json}
	\item \texttt{server/app.py}
	\item \texttt{client/src/components/TokenDashboardModal.jsx}
\end{itemize}

\begin{notebox}[Neden Bu Kapsam?]
Extension akisinda komut, dogrulama ve istek gonderimi; token akisinda paket, checkout ve webhook adimlari birbiriyle bagimli oldugu icin dokumanin odagi bu dosyalardir.
\end{notebox}

\section{VS Code Extension Adimlari}

\subsection{Adim 1 -- Uzanti kontratinin tanimlanmasi}

\texttt{vscode-extension/package.json} uzerinde extension kontrati netlestirilmistir:
\begin{itemize}[leftmargin=1.5em]
	\item Iki komut tanimlanmistir: \texttt{codeAlchemist.setApiKey} ve \texttt{codeAlchemist.ask}.
	\item Aktivasyon tetikleyicileri bu komutlara baglanmistir.
	\item Ayarlar altinda \texttt{codeAlchemist.apiKey} ve \texttt{codeAlchemist.endpoint} konfigurasyon anahtarlari acilmistir.
\end{itemize}

Bu adimla extension'in calisma kosullari, koddan once manifest seviyesinde garanti altina alinmistir.

\subsection{Adim 2 -- API anahtari yonetimi}

\texttt{extension.ts} icinde \texttt{CodeAlchemist: Set API Key} komutu su davranisi saglar:
\begin{itemize}[leftmargin=1.5em]
	\item Kullaniciya gizli giris kutusu acar.
	\item API anahtarini \texttt{VS Code SecretStorage} icinde saklar.
	\item Eski surumlerde settings icinde saklanmis \texttt{codeAlchemist.apiKey} degeri varsa ilk kullanimda otomatik migrate eder.
	\item Basarili kayitta kullaniciya bilgi mesaji gosterir.
\end{itemize}

Bu tasarimla anahtar hem hard-code edilmez hem de plain settings yerine guvenli kasada tutulur.

\subsection{Adim 3 -- Soru gonderim akisi (Ask komutu)}

\texttt{codeAlchemist.ask} komutu calistiginda asagidaki siralama izlenir:
\begin{enumerate}[leftmargin=1.5em]
	\item Ayarlardan API key ve endpoint okunur.
	\item Eksik degilse aktif editor secimi alinir (varsa secili kod da mesaja eklenir).
	\item Kullanicidan soru metni alinir.
	\item \texttt{POST} istegi ile endpoint'e JSON govde gonderilir: \texttt{\{question, code\}}.
	\item Sunucu \texttt{text/event-stream} donuyorsa cevap parca parca (SSE) canli olarak Output kanalina akar.
	\item Stream yoksa JSON fallback ile toplu cevap yazdirilir.
\end{enumerate}

\subsection{Adim 4 -- Sunucu tarafi extension endpoint'i}

\texttt{server/app.py} icindeki \texttt{/v1/ask} endpoint'i extension icin dis entegrasyon katmanidir.

Beklenen sozlesme:
\begin{itemize}[leftmargin=1.5em]
	\item Header: \texttt{X-API-Key}
	\item Body: \texttt{question} ve/veya \texttt{code}
\end{itemize}

Ana davranis:
\begin{itemize}[leftmargin=1.5em]
	\item API key dogrulanir (aktif anahtar kontrolu).
	\item Anahtarin son kullanim zamani guncellenir.
	\item Gerekli giris varsa Gemini cevap uretecisi tetiklenir.
	\item Istek \texttt{text/event-stream} bekliyorsa stream, aksi halde toplu JSON cevap donulur.
\end{itemize}

\subsection{Adim 5 -- API key yasam dongusu}

Kullanicinin extension icin anahtar uretip yonetebilmesi su endpoint seti ile saglanmistir:
\begin{itemize}[leftmargin=1.5em]
	\item \texttt{GET /api/keys}: aktif anahtarlarin listelenmesi
	\item \texttt{POST /api/keys}: yeni \texttt{ca-...} formatinda anahtar olusturma
	\item \texttt{DELETE /api/keys/<id>}: soft revoke (is_active=false)
\end{itemize}

\section{Token Satin Alma Mekanizmasi}

\subsection{Mimari ozet}

Token odeme sistemi 3 katmanda calisir:
\begin{enumerate}[leftmargin=1.5em]
	\item \textbf{Paket ve bakiye katmani}: \texttt{TokenPackage}, \texttt{TokenBalance}, \texttt{TokenTransaction}, \texttt{TokenPurchase} modelleri.
	\item \textbf{Odeme orkestrasyonu}: \texttt{/api/billing/checkout-session} ve \texttt{/api/billing/webhook}.
	\item \textbf{UI gorunurlugu}: \texttt{TokenDashboardModal.jsx} ile paket secimi, gecmis, bakiye panelleri.
\end{enumerate}

\subsection{Paketlerin olusturulmasi ve servis edilmesi}

Sistem acilisinda \texttt{ensure\_default\_token\_packages()} fonksiyonu aktif paketleri seed eder.
Frontend, paketleri \texttt{GET /api/billing/packages} ile ceker.

\begin{table}[h]
\centering
\begin{tabularx}{\textwidth}{@{}lrrrX@{}}
\toprule
Paket & Token & Fiyat (USD) & Bonus \% & Not \\
\midrule
Starter & 500 & 5 & 0 & Hafif deneme akisi \\
Pro Pack & 2000 & 15 & 5 & Varsayilan one cikan paket \\
Heavy User Bundle & 8000 & 49 & 10 & Yogun kullanim \\
Studio Upgrade & 20000 & 99 & 15 & Yuksek hacimli senaryo \\
\bottomrule
\end{tabularx}
\caption{Token paketleri (varsayilan kurulum)}
\end{table}

\subsection{Checkout oturumu olusturma}

\texttt{POST /api/billing/checkout-session} akisinda:
\begin{enumerate}[leftmargin=1.5em]
	\item JWT ile kullanici dogrulanir.
	\item \texttt{package\_id} once DB'de, bulunamazsa varsayilan paket listesinde eslestirilir.
	\item Bonus token hesaplanir: toplam token = temel token + bonus.
	\item Stripe Checkout Session olusturulur.
	\item Veritabanina \texttt{pending} durumunda \texttt{TokenPurchase} kaydi yazilir.
	\item Frontend'e yonlendirme icin \texttt{checkout\_url} donulur.
\end{enumerate}

\subsection{Webhook ile token yukleme}

\texttt{POST /api/billing/webhook} endpoint'i Stripe'dan gelen \texttt{checkout.session.completed} olayini dinler.

Tamamlanma adimlari:
\begin{itemize}[leftmargin=1.5em]
	\item Imza dogrulamasi yapilir (webhook secret).
	\item Ilgili satin alma kaydi bulunur; yoksa metadata ile fallback kaydi uretilir.
	\item Idempotency kontrolu: satin alma zaten \texttt{completed} ise tekrar kredi verilmez.
	\item \texttt{TokenBalance} guncellenir.
	\item \texttt{TokenTransaction(type='purchase')} kaydi eklenir.
	\item Satin alma kaydi \texttt{completed} ve \texttt{completed\_at} alanlariyla kapanir.
\end{itemize}

Bu yapi, cift webhook veya yeniden gonderim durumlarinda bakiye sapmasini engellemek icin kritik onemdedir.

\section{Hafta 9.1 Uygulanan Iyilestirmeler}

\subsection{1) Extension -- SecretStorage gecisi}

Uygulanan degisiklikler:
\begin{itemize}[leftmargin=1.5em]
	\item API key saklama noktasi settings'ten \texttt{context.secrets} tabanli SecretStorage'a alindi.
	\item Legacy ayar degeri tespit edilirse otomatik tasima (migration) eklendi.
	\item Bos anahtar girildiginde SecretStorage kaydi temizleniyor.
\end{itemize}

Teknik kazanim: API key'in editor ayarlarinda acik metin olarak kalma riski azaltildi.

\subsection{2) Extension -- SSE canli cikti}

Uygulanan degisiklikler:
\begin{itemize}[leftmargin=1.5em]
	\item \texttt{Accept: text/event-stream, application/json} header'i eklendi.
	\item SSE \texttt{data: ...} satirlari stream reader ile parse edilip Output kanalina anlik yazdiriliyor.
	\item \texttt{[DONE]} sinyali ve stream sonu durumlari guvenli sekilde ele aliniyor.
	\item Stream disi cevaplarda JSON/text fallback korunuyor.
\end{itemize}

Teknik kazanim: Kullanici tek parca sonuc beklemek yerine cevabi token/parca bazinda canli goruyor.

\subsection{3) Billing -- Webhook idempotency sertlestirme}

Uygulanan degisiklikler:
\begin{itemize}[leftmargin=1.5em]
	\item \texttt{checkout.session.completed} olayinda session id yoksa islem erken sonlandiriliyor.
	\item Satin alma kaydi atomik claim mantigi ile \texttt{pending -> processing} gecisine zorlandi.
	\item Ayni \texttt{reference\_id} icin mevcut \texttt{purchase} transaction kaydi varsa kredi tekrar yuklenmiyor.
	\item Gecis tamamlaniyorsa satin alma \texttt{completed} olarak kapanip audit izi korunuyor.
\end{itemize}

Teknik kazanim: cift webhook, tekrar gonderim ve yaris kosullarinda cift bakiye yazimi riski belirgin sekilde azaldi.

\subsection{4) Billing -- Checkout sonrasi otomatik bakiye yenileme}

Uygulanan degisiklikler:
\begin{itemize}[leftmargin=1.5em]
	\item Frontend URL query'sinde \texttt{payment=success} veya \texttt{billing=success} algilandiginda otomatik yenileme akisi tetikleniyor.
	\item \texttt{/api/tokens/usage} cagrisi ile bakiye aliniyor ve kullanici state/localStorage guncelleniyor.
	\item \texttt{/api/billing/usage} verisi yenileniyor ve Token Dashboard otomatik aciliyor.
	\item URL query parametreleri islendikten sonra temizleniyor (tekrar tetiklemeyi onlemek icin).
\end{itemize}

Teknik kazanim: odeme donusunde kullanici manuel yenileme yapmadan guncel bakiyeyi goruyor.

\subsection{Frontend satin alma deneyimi}

\texttt{TokenDashboardModal.jsx} tarafinda kullanici akisi:
\begin{enumerate}[leftmargin=1.5em]
	\item Modal acildiginda paralel 3 istek atilir:
	\texttt{/api/tokens/usage}, \texttt{/api/billing/packages}, \texttt{/api/billing/config}.
	\item Paket seciminde \texttt{handleCheckout} ile \texttt{/api/billing/checkout-session} cagrilir.
	\item Donen \texttt{checkout\_url} ile tarayici Stripe sayfasina yonlendirilir.
	\item Odeme sonrasi webhook tokenlari yukler; kullanici tekrar modal actiginda bakiye ve islem gecmisi guncel gorunur.
\end{enumerate}

\subsection{Token kullanim gorunurlugu}

\texttt{GET /api/tokens/usage} endpoint'i su verileri dondurur:
\begin{itemize}[leftmargin=1.5em]
	\item Guncel bakiye
	\item Toplam harcanan token
	\item Son N islem (alim ve kullanim kayitlari)
\end{itemize}

Bu veri, UI'da kullanicinin ne kadar token harcadigini izleyebilmesi ve satin alma kararini veriye dayali alabilmesi icin kullanilir.

\section{Uctan Uca Surec Ozetleri}

\subsection{Flow A -- VS Code uzerinden soru-cevap}

\begin{enumerate}[leftmargin=1.5em]
	\item Kullanici extension ayarindan API key ve endpoint tanimlar.
	\item Editor icinde kod secip \texttt{Ask} komutunu tetikler.
	\item Extension, \texttt{X-API-Key} ile \texttt{/v1/ask} endpoint'ine istek yollar.
	\item Sunucu model yanitini uretir ve extension Output kanalinda sonucu gosterir.
\end{enumerate}

\subsection{Canli Dogrulama Notu -- Extension Testi}

9 Nisan 2026 tarihinde Extension Development Host uzerinde canli test gerceklestirilmistir.
Testte \texttt{SELAM} sorusu ile \texttt{POST /v1/ask} cagrisi atilmis ve asagidaki beklenen ciktilar alinmistir:

\begin{itemize}[leftmargin=1.5em]
	\item \texttt{Status: 200 OK}
	\item Sistem fallback mesaji: \texttt{Previous model failed, trying ...}
	\item Fallback sonrasi modelden tam metin yanitinin Output kanalinda gorunmesi
\end{itemize}

Bu sonuc, extension tarafinda hem istek-gonderim hattinin hem de stream/fallback goruntuleme akisinin calistigini dogrular.

\subsection{Flow B -- Token satin alma ve bakiye artisi}

\begin{enumerate}[leftmargin=1.5em]
	\item Kullanici Token Dashboard'dan paket secer.
	\item Frontend checkout-session olusturur ve Stripe'a yonlendirir.
	\item Stripe, odeme tamamlandiginda webhook gonderir.
	\item Sunucu satin almayi tamamlar, cuzdana token ekler, transaction yazar.
	\item Kullanici panelde yeni bakiyeyi gorur.
\end{enumerate}

\section{Guvenlik ve Dayaniklilik Notlari}

\begin{itemize}[leftmargin=1.5em]
	\item API key dogrulamasi extension endpoint'inde zorunludur; gecersiz anahtarlar reddedilir.
	\item Satin alma akisinda hem Stripe imza dogrulamasi hem de idempotency kontrolu bulunur.
	\item Bakiye artisi sadece webhook tarafinda \texttt{completed} event ile yapildigi icin istemci manipule edilemez.
	\item Satin alma kaydinin \texttt{pending -> completed} gecisi denetlenebilir bir audit izi olusturur.
\end{itemize}

\section{Hafta 9.2 Baglam Koruma Modulu}

Bu guncellemenin hedefi, kullanicinin yeni sohbet actiginda onceki teknik baglami tekrar tekrar anlatma ihtiyacini azaltmak ve token israfini onlemektir.

\subsection{Problem Tanimi}

Ornek senaryo:
\begin{itemize}[leftmargin=1.5em]
	\item Kullanici bir website tasarlarken birden fazla sohbet aciyor.
	\item Yeni sohbette model, onceki kararlar (stack, UI tarzi, kisitlar) bilinmedigi icin ayni baglam tekrar yaziliyor.
	\item Bu tekrar, gecikmeyi ve token tuketimini arttiriyor.
\end{itemize}

\subsection{Uygulanan Teknik Cozum}

\textbf{1) Otomatik devam modu (backend default-on):}
\begin{itemize}[leftmargin=1.5em]
	\item \texttt{server/app.py} icine \texttt{\_resolve\_include\_previous\_modules(...)} yardimci fonksiyonu eklendi.
	\item Authenticated chat isteklerinde, kullanici acikca kapatmadiysa onceki baglam otomatik dahil ediliyor.
	\item \texttt{no\_save=true} isteklerinde memory carryover kapali kalmaya devam ediyor.
\end{itemize}

\textbf{2) Hibrit retrieval ile baglam secimi:}
\begin{itemize}[leftmargin=1.5em]
	\item \texttt{server/utils/memory\_utils.py} icinde lexical + semantic (embedding) + recency skorlamasi birlikte kullaniliyor.
	\item Embedding ulasilamazsa lexical fallback devam ediyor; sistem hard-fail olmuyor.
	\item Ciktiya \texttt{retrieval\_mode} ve \texttt{focus\_module} metadata'si eklenerek izlenebilirlik saglaniyor.
\end{itemize}

\textbf{3) Website tasarimina ozel memory kurallari:}
\begin{itemize}[leftmargin=1.5em]
	\item Yeni modul anahtarlari eklendi: \texttt{project\_scope}, \texttt{ui\_style}, \texttt{tech\_stack}.
	\item Ornek sinyaller: \texttt{website}, \texttt{landing page}, \texttt{responsive}, \texttt{tailwind}, \texttt{react}.
	\item Bu sayede yeni sohbette "siteyi responsive devam edelim" gibi kisa istemlerde onceki kararlar korunuyor.
\end{itemize}

\textbf{4) Frontend varsayilan davranis:}
\begin{itemize}[leftmargin=1.5em]
	\item \texttt{client/src/App.jsx} icinde \texttt{includePreviousModules} ilk kullanimda varsayilan \texttt{true} yapildi.
	\item Kullanici daha once secim yaptiysa localStorage degeri korunuyor.
\end{itemize}

\subsection{Dogrulama ve Test}

\begin{itemize}[leftmargin=1.5em]
	\item Unit test suiti: \texttt{python -m unittest server.tests.test\_memory\_utils}
	\item Sonuc: Tum testler \texttt{OK} (extract, budget/limit, summary fallback, semantic tercih, website memory extraction).
\end{itemize}

\subsection{Beklenen Urun Etkisi}

\begin{itemize}[leftmargin=1.5em]
	\item Yeni sohbetlerde tekrar baglam anlatimi azalir.
	\item Prompt uzunlugu kisalir, token tuketimi duser.
	\item Kullanici deneyimi "kaldigim yerden devam" seviyesine gelir.
\end{itemize}

\subsection{Urun Seviyesi Memory Mimarisine Gecis}

Bu modul, "tam gecmis enjekte et" yaklasimindan uzaklasip 3 katmanli bellek modeline yaklastirilmistir:

\begin{enumerate}[leftmargin=1.5em]
	\item \textbf{Working memory}: aktif sohbetin anlik baglami.
	\item \textbf{Episodic memory}: onceki oturum ozetleri ve kararlar (retrieval ile secilir).
	\item \textbf{Semantic memory}: daha kalici tercih/kisit sinyalleri (module bazli kararlar).
\end{enumerate}

Bu geciste temel ilke su sekilde uygulanmistir:
\begin{itemize}[leftmargin=1.5em]
	\item Full context carryover yerine maksimum 5 satirlik \texttt{[Memory Capsule]} kullanimi.
	\item Confidence tabanli filtreleme ile dusuk alaka skorlu belleklerin prompta girmemesi.
	\item "Sifirdan basla" niyeti algilandiginda carryover'in tamamen atlanmasi (decision lock-in onleme).
	\item Varsayilan olarak sadece son oturum ozetinin retrieval icinde dikkate alinmasi.
\end{itemize}

\subsection{Frontend Test Senaryolari}

Baglam koruma davranisinin frontend'de dogrulanmasi icin asagidaki senaryolar kullanilmistir:
\begin{enumerate}[leftmargin=1.5em]
	\item Browser localStorage icinde \texttt{includePreviousModules} degerinin ilk yuklemede \texttt{true} geldigini kontrol et.
	\item Yeni sohbet acip onceki proje kararlarinin \texttt{[Memory Capsule]} olarak istege eklendigini gozlemle.
	\item \texttt{no\_save=true} veya "sifirdan basla" akisi ile carryover'in kapandigini dogrula.
	\item Ayni konu icin tekrar soru gonderip cevaplarda tekrar baglam yazimi azaliyorsa memory injection'in calistigini kontrol et.
	\item Network panelde istek payload'i icinde soru + secili kod + memory metadata akisini incele.
\end{enumerate}

Bu testler, kullanicinin her yeni sohbette baglami yeniden anlatmadan devam edebilmesini ve frontend tarafinda memory davranisinin gorunur olmasini saglar.

\section{Hafta 9.3 Frontend Debug Paneli}

Memory/carryover davranisinin frontend'de gorunur olmasi ve test edilebilmesi icin React state tabanli bir debug paneli eklenmistir.

\subsection{Teknik Uygulama}

\textbf{1) State tanimlamasi:}
\begin{itemize}[leftmargin=1.5em]
	\item \texttt{client/src/App.jsx} icinde debug state eklendi: \texttt{const [debug, setDebug] = useState(null)}
	\item State null ile baslatilarak panel baslangicta gorunmez.
\end{itemize}

\textbf{2) API response yakalama:}
\begin{itemize}[leftmargin=1.5em]
	\item Stream bitirme sinyali (\texttt{data.done}) alindiktan sonra debug verileri yakalanir.
	\item Yakalanan veri: \texttt{carryover}, \texttt{memory\_capsules} array'i ve \texttt{capsuleCount}.
	\item Backend'den gelen response'ta bu alanlar zorunlu olusturulur.
\end{itemize}

\texttt{setDebug(\{}\texttt{carryover: data.carryover, capsuleCount: data.memory\_capsules?.length || 0, capsules: data.memory\_capsules || []\}\texttt{);}

\textbf{3) UI Render:}
\begin{itemize}[leftmargin=1.5em]
	\item Debug kutusu ekranin sag alt kosesinde fixed position ile yerlestirilmis.
	\item Arka plan: opasite ile \texttt{\#111} (siyah).
	\item Font: monospace, ufak font size ile read-only gorunumu.
	\item Border: hafif gri renkte (\texttt{\#555}).
	\item Z-index: 999 (Dialog'un ustunde goruntu).
\end{itemize}

\subsection{Goruntulenen Bilgiler}

Debug paneli asagidaki verileri sunmustur:

\begin{table}[h]
\centering
\begin{tabularx}{\textwidth}{@{}lX@{}}
\toprule
Alan & Anlam \\
\midrule
\texttt{Carryover: true/false} & Onceki baglam dahil edilip edilmediginin simulasyonu (yesil=true, kirmizi=false) \\
\texttt{Capsules: N} & Aktif memory capsule sayisi \\
\texttt{Capsule Listeleri} & Her capsule'in \texttt{summary} veya \texttt{title} ozeti \\
\bottomrule
\end{tabularx}
\caption{Frontend debug paneli goruntulenen veri alanlar}
\end{table}

\subsection{Test Adimlari}

\begin{enumerate}[leftmargin=1.5em]
	\item \textbf{Normal sohbet}: Soru gonder, cevap al. Sag alt kosede debug paneli gorulu olmali.
	\item \textbf{Carryover kontrol}: 
	\begin{itemize}[leftmargin=1.5em]
		\item Carryover: true ise yeni sohbette onceki veri eklenecek.
		\item Carryover: false ise yeni baslangic oldugu anlamina gelir.
	\end{itemize}
	\item \textbf{Capsule sayisi}: Artis gozlemle. Yeni sohbet her istekte daha fazla capsule ekleyebilir.
	\item \textbf{no\_save true durum}: Eger \texttt{no\_save=true} ile istek yapilirsa carryover false olmali.
	\item \textbf{Network dogrulama}: Browser DevTools Network tab'inda request payload'inde:
	\begin{itemize}[leftmargin=1.5em]
		\item \texttt{include\_previous\_modules: true/false}
		\item Memory capsule enjeksiyonu gorunmeli.
	\end{itemize}
\end{enumerate}

\subsection{Backend Entegrasyonu}

Debug panelinın calisabilmesi icin \texttt{/api/ask} ve \texttt{/api/blend} endpoint'leri response'ta su alanlari donmeli:

\begin{itemize}[leftmargin=1.5em]
	\item \texttt{carryover}: boolean - onceki baglam kullanilip kullanilmadigi
	\item \texttt{memory\_capsules}: array - secili memory capsule'ler listesi
	\item Her capsule'de en az \texttt{summary} veya \texttt{title} alan olmali
\end{itemize}

\subsection{Kullanici Rahatyeti}

Debug panel:
\begin{itemize}[leftmargin=1.5em]
	\item Tamamen opsiyoneldir (production'da gizlenebilir).
	\item Sabit konumda (sag alt) kalir, skrol edilmez.
	\item Minimal CSS ile performans etkilemez.
	\item Hicbir kullanici etkisini almaz (read-only).
\end{itemize}


\section{Hafta 9.4 Agent Modu -- Entegrasyon ve Optimizasyon}

Bu bolum, CodeAlchemist platformuna eklenen Agent Modu ozelliginin teknik mimarisini, entegrasyon adimlarini ve ardindan uygulanan optimizasyonlari kapsamaktadir.

\subsection{Agent Modu Nedir?}

Agent Modu, klasik tek donus (request-response) sohbet modelinin otesine gecerek LLM'nin \textit{arac kullanabilmesini} (tool use) ve \textit{cok adimli gorev yurutebilmesini} saglayan bir calisma modudur.

Standart moddaki farki:
\begin{itemize}[leftmargin=1.5em]
    \item \textbf{Standart Mod}: Kullanici mesaj gonderir $\rightarrow$ model tek bir yanit uretir $\rightarrow$ akis biter.
    \item \textbf{Agent Modu}: Kullanici mesaj gonderir $\rightarrow$ model hedefi analiz eder $\rightarrow$ gerekli araclari (web arama, kod analizi, dosya okuma vb.) cagirip sonuclari degerlendirerek nihai yaniti olusturur.
\end{itemize}

\begin{notebox}[Agent Modu Kullanim Ornekleri]
``Bugunku hava durumu nedir?'' $\rightarrow$ web arama araci kullanilir.\\
``Bu Python kodundaki hatayi bul ve duzelt.'' $\rightarrow$ kod analiz araci devreye girer.\\
``Bu projedeki dosyalara gore mimari ozet cikar.'' $\rightarrow$ proje dosya okuma araci kullanilir.
\end{notebox}

\subsection{Mimari Yapi}

Agent Modu, hibrit bir mimaride hayata gecirilmistir:

\begin{table}[h]
\centering
\begin{tabularx}{\textwidth}{@{}lX@{}}
\toprule
Katman & Aciklama \\
\midrule
\texttt{server/services/agent\_runtime.py} & Legacy Flask runtime; \texttt{run\_agent\_turn} giris noktasi \\
\texttt{server/backend/runtime/context.py} & Yeni FastAPI tabanli unified runtime; \texttt{ContextAssembler.assemble} \\
\texttt{server/services/agent\_bridge.py} & Flask $\leftrightarrow$ AgentRuntime koprusu; SSE aktarimini yonetir \\
\texttt{server/app.py} & \texttt{/api/ask} endpoint'i; \texttt{agent\_mode} parametresini alir \\
\texttt{client/src/App.jsx} & \texttt{agentModeEnabled} state ve toggle UI \\
\texttt{client/src/components/ChatInterface.jsx} & \texttt{AgentTracePanel} -- gercek zamanli arac izi goruntuleyici \\
\bottomrule
\end{tabularx}
\caption{Agent Modu mimarisi -- katmanlar ve dosyalar}
\end{table}

\subsection{Entegrasyon Adimlari}

\textbf{1) Backend -- Agent Bridge:}
\begin{itemize}[leftmargin=1.5em]
    \item \texttt{/api/ask} endpoint'ine \texttt{agent\_mode} boolean parametresi eklendi.
    \item \texttt{agent\_mode=True} geldiginde \texttt{stream\_agent\_bridge()} cagrilarak istek \texttt{AgentRuntime}'a yonlendirilir.
    \item \texttt{AgentRuntime}, mevcut arac kayit defterini (\texttt{tool\_registry}) kullanarak LLM cevrimleri baslatir.
    \item Her arac cagrisi ve LLM adimi SSE (Server-Sent Events) uzerinden \\ anlık olarak frontend'e iletilir.
\end{itemize}

\textbf{2) Sistem Promptu:}
\begin{itemize}[leftmargin=1.5em]
    \item Her iki runtime icin farkli \texttt{build\_system\_prompt} fonksiyonlari mevcuttu.
    \item Her ikisine de Agent Modu karar mantigi talimatlari eklendi (bkz. Bolum 9.4.5).
\end{itemize}

\textbf{3) Frontend -- Toggle ve Trace:}
\begin{itemize}[leftmargin=1.5em]
    \item \texttt{App.jsx} icinde \texttt{agentModeEnabled} state tanimlanmis; kullanici bu dugmeye basarak Agent Modunu acip kapayabiliyor.
    \item \texttt{agent\_mode: Boolean(agentModeEnabled)} olarak her POST istegine ekleniyor.
    \item \texttt{ChatInterface.jsx} icindeki \texttt{AgentTracePanel} bileseni, SSE olaylarini dinleyerek arac cagrilari, adim sayisi ve sonuclari gercek zamanli gosteriyor.
\end{itemize}

\textbf{4) SSE Streaming Protokolu:}

Agent Modu aktifken sunucu su tipte olaylar yayinlar:
\begin{itemize}[leftmargin=1.5em]
    \item \texttt{type: "message"} -- LLM'den gelen metin parcalari.
    \item \texttt{type: "tool\_call"} -- Bir aracin cagrildigini bildiren olay.
    \item \texttt{type: "tool\_result"} -- Aracin donduhu sonuc.
    \item \texttt{type: "done"} -- Akis bitmis; \texttt{trace} ve \texttt{changed\_files} listeleri eklenmis.
\end{itemize}

\texttt{agent\_bridge.py} bu olaylari legacy UI formatina donusturur (\texttt{chunk} anahtari).

\subsection{Desteklenen Modeller}

Agent Modu, arac cagrisini (function calling / tool use) destekleyen modeller ile calisir:
\begin{itemize}[leftmargin=1.5em]
    \item \textbf{Gemini}: \texttt{gemini-2.5-flash}, \texttt{gemini-2.5-pro}
    \item \textbf{GPT}: \texttt{gpt-4o}, \texttt{gpt-4o-mini}
    \item \textbf{Claude}: \texttt{claude-3-5-sonnet}, \texttt{claude-opus-4-5}
\end{itemize}

Desteklenmeyen modellerde Agent Modu secilirse kullaniciya aciklayici bir hata mesaji dondurulur.

\subsection{Optimizasyon Gereksinimi -- Problem Tanimi}

Agent Modu entegrasyonunun ardindan iki kritik sorun tespit edildi:
\begin{itemize}[leftmargin=1.5em]
    \item \textbf{Gereksiz Agent Tetiklemesi}: ``selam'', ``tesekkurler'' gibi basit sohbet mesajlari da tum Agent arac cagri altyapisini devreye sokuyordu; bu durum gecikmeye ve token israfina yol aciyordu.
    \item \textbf{Kayip Sohbetler}: Agent Modu acikken baslanan yeni sohbetler, sistemin ``Auto-Discovery'' mekanizmasi tarafindan kullanicinin son projesine otomatik olarak baglaniyordu. Bu nedenle bu sohbetler sol menudeki Conversations listesinde gorunmuyordu.
\end{itemize}

\subsection{Cozum 1 -- Yonlendirme Prompt Enjeksiyonu}

Her iki agent runtime'ina da yonlendirme talimatlari eklendi:
\begin{itemize}[leftmargin=1.5em]
    \item \textbf{Legacy Runtime}: \texttt{server/services/agent\_runtime.py} icindeki \texttt{build\_agent\_system\_prompt} fonksiyonuna Turkce yonlendirme blogu eklendi.
    \item \textbf{Yeni Runtime}: \texttt{server/backend/runtime/context.py} icindeki \texttt{\_build\_system\_prompt} fonksiyonuna ayni yonlendirme blogu eklendi.
\end{itemize}

Eklenen talimat blogu modele iki kategori tanimladi:
\begin{itemize}[leftmargin=1.5em]
    \item \textbf{AGENT MODUNU KULLAN}: Web aramasi, kod yazma, dosya isleme, cok adimli gorevler, gercek zamanli veri.
    \item \textbf{NORMAL YANIT VER}: Selamlama, tesekkur, kisa tanim sorulari, evet/hayir sorulari.
\end{itemize}

Enjeksiyon, modelin son gordugundan etkilenmesi nedeniyle prompt sonuna yapildi (append stratejisi).

\subsection{Cozum 2 -- Heuristik Koruma Katmani (Guardrail)}

LLM'nin her zaman prompt talimatlarina uymasi beklenemeyecegi icin \texttt{server/app.py} icine \texttt{/api/ask} endpoint'inde hafif bir ontest mekanizmasi eklendi.

Gelen mesaj 5 kelimeden kisa ise su kontroller uygulanir:
\begin{itemize}[leftmargin=1.5em]
    \item Mesaj bilinen sohbet kelimeleri iceriyorsa (selam, merhaba, tesekkur, hi vb.) Agent Modu iptal edilir.
    \item Mesaj ``nedir'' veya ``kimdir'' icerip 3 kelimeden kisaysa Agent Modu iptal edilir.
\end{itemize}

Agent Modu iptal edildiginde kullaniciya dogrudan \texttt{routing\_reason} aciklamasi dondurulur.

\textbf{Test Sonuclari} (\texttt{server/test\_routing.py}):

\begin{table}[h]
\centering
\begin{tabularx}{\textwidth}{@{}lccX@{}}
\toprule
Girdi & Beklenen & Sonuc & Aciklama \\
\midrule
selam & Normal & \checkmark PASS & Sohbet kelimesi \\
hava durumu nedir & Agent & \checkmark PASS & 3+ kelime, arac gerekli \\
tesekkurler & Normal & \checkmark PASS & Sohbet kelimesi \\
python ile fibonacci yaz & Agent & \checkmark PASS & Kod gorevi \\
api nedir & Normal & \checkmark PASS & Kisa tanim, 2 kelime \\
su log dosyasini analiz et & Agent & \checkmark PASS & Analiz gorevi \\
\bottomrule
\end{tabularx}
\caption{Routing heuristic test sonuclari (6/6 basarili)}
\end{table}

\subsection{Cozum 3 -- Otomatik Proje Baglama (Auto-Link) Kaldirildi}

Onceki surumde \texttt{\_resolve\_agent\_project} fonksiyonunda bir ``Auto-Discovery'' mekanizmasi mevcuttu:
\begin{itemize}[leftmargin=1.5em]
    \item Agent Modu acikken baslanan her yeni sohbet, kullanicinin en son aktif projesine otomatik baglaniyordu.
    \item Proje bagli sohbetler \texttt{GET /api/conversations} endpoint'inde \texttt{project\_id IS NULL} filtresi nedeniyle Conversations listesinde gorunmuyordu.
\end{itemize}

Uygulanan degisiklikler:
\begin{itemize}[leftmargin=1.5em]
    \item \texttt{app.py} icindeki auto-fallback kod blogu tamamen kaldirildi; yorumlanmis sekilde log ucak birakildi.
    \item \texttt{server/models.py} icindeki \texttt{Conversation} modeline \texttt{is\_auto\_linked} Boolean alani eklendi. Bu alan, sistem tarafindan baglanan sohbetleri kullanici el ile bagladigi sohbetlerden ayirt etmeyi saglar.
\end{itemize}

\begin{notebox}[Duzeltilen Davranis]
Onceki hali: \texttt{agent\_mode=True} $\rightarrow$ son projeye otomatik bagla (sohbet gizlenir).\newline
Yeni hali: \texttt{agent\_mode=True} $\rightarrow$ Conversations listesinde gorunur; kullanici isterse manuel projeye ekler.
\end{notebox}

\subsection{Cozum 4 -- Veritabani Migrasyonu ve Gecmis Sohbet Kurtarma}

Ayni oturumda onceden kaybolan sohbetleri geri getirmek icin \texttt{server/migrate\_auto\_link.py} migration scrip'i gelistirildi ve calistirildi.

Migration adimlari:
\begin{enumerate}[leftmargin=1.5em]
    \item \texttt{conversation} tablosuna \texttt{is\_auto\_linked} kolonu eklendi (mevcut kayitlar False ile dolduruldu).
    \item Son 6 saat icinde olusturulan, \texttt{project\_id IS NOT NULL} olan sohbetler heuristik olarak auto-linked kabul edildi ve isaretlendi.
    \item Bu sohbetlerin \texttt{project\_id} degeri NULL yapilarak Conversations listesine geri tasinmasi saglandi.
\end{enumerate}

\textbf{Migration Ciktisi:}
\begin{itemize}[leftmargin=1.5em]
    \item \texttt{is\_auto\_linked} kolonu basariyla eklendi.
    \item 19 sohbet yeniden ``genel'' olarak isaretlen.
    \item 19 sohhet sol menudeki Conversations listesine geri kazandirildi.
\end{itemize}

\subsection{Etkilenen Dosyalar}

\begin{table}[h]
\centering
\begin{tabularx}{\textwidth}{@{}lX@{}}
\toprule
Dosya & Degisiklik \\
\midrule
\texttt{server/app.py} & Guardrail blogu eklendi; auto-fallback kaldirildi \\
\texttt{server/models.py} & \texttt{is\_auto\_linked} alani eklendi \\
\texttt{server/services/agent\_runtime.py} & Routing talimatlari prompt sonuna eklendi \\
\texttt{server/backend/runtime/context.py} & Routing talimatlari prompt sonuna eklendi \\
\texttt{server/test\_routing.py} & Tam test suiti guncellendi (6/6 basarili) \\
\texttt{server/migrate\_auto\_link.py} & Yeni migration scrip'i olusturuldu ve calistirildi \\
\bottomrule
\end{tabularx}
\caption{9.4 guncellemesinde etkilenen dosyalar}
\end{table}

\subsection{Mimari Karar: Neden Auto-Link Kaldirildi?}

Bu karar, buyuk AI platformlarinin (ChatGPT, Claude, Gemini) proje yonetim modeliyle hizalanmak icin alinmistir:
\begin{itemize}[leftmargin=1.5em]
    \item Kullanici nereye kaydedeceğine kendisi karar verir.
    \item Sistem asla otomatik baglamaz.
    \item Temiz ayirim (general chat vs. project workspace) daha iyi UX saglar.
    \item Kullanici isterse sohbeti sonradan projeye ekleyebilir; bu dogal ve anlasilir bir akistir.
\end{itemize}

\section{Hafta 9.5 Extension-Docker Full Entegrasyon ve UTF-8 Hata Giderimi}

Bu bolumde, VS Code extension ile Docker uzerinde calisan backend arasindaki uctan uca entegrasyon testi ve ayni surecte giderilen JWT/encoding kaynakli hata sinifi ozetlenmektedir.

\subsection{Problem Ozeti}

Sahada gozlenen iki kritik problem su sekildeydi:
\begin{itemize}[leftmargin=1.5em]
	\item \textbf{JWT decode kaynakli UTF-8 hatasi}: Gecersiz veya bozuk \texttt{Authorization} header degerleri JWT kutuphanesine ulastiginda \texttt{Invalid header string} ve \texttt{utf-8 codec} hatalari uretebiliyordu.
	\item \textbf{Auth akisi karisimi riski}: API key ve JWT token akislari farkli endpoint kontratlarina sahip oldugu halde ayni header kanalina dusen gecersiz degerler 500 sinifi davranisa yaklasabiliyordu.
\end{itemize}

\subsection{Uygulanan Cozum}

\textbf{1) Erken asama header sanitization (before\_request):}
\begin{itemize}[leftmargin=1.5em]
	\item \texttt{server/app.py} icine \texttt{sanitize\_auth\_headers()} hook'u eklendi.
	\item \texttt{Bearer} tipinde gelen degerler icin JWT segment yapisi (3 parcali token) dogrulanir hale getirildi.
	\item Beklenmeyen formatlar veya encoding supheleri tespit edildiginde \texttt{HTTP\_AUTHORIZATION} temizlenerek JWT katmanina bozuk veri gecisi engellendi.
\end{itemize}

\textbf{2) JWT ve API key akisinin ayrik tutulmasi:}
\begin{itemize}[leftmargin=1.5em]
	\item Extension entegrasyonunda \texttt{/v1/ask} icin \texttt{X-API-Key} zorunlu kaldi.
	\item Kullaniciya bagli endpoint'lerde \texttt{Authorization: Bearer <jwt>} akisi korunarak iki kimlikleme tipi net sekilde ayrildi.
\end{itemize}

\textbf{3) Locale-guvenli sayisal serilestirme:}
\begin{itemize}[leftmargin=1.5em]
	\item \texttt{progress\_percent} alani icin acik \texttt{float()} donusumu uygulanarak JSON ciktilarinda ondalik ayiraci tutarlilastirildi.
	\item Boylece TR locale ortamlarinda \texttt{1,1} yerine RFC uyumlu \texttt{1.1} cikisi garanti altina alindi.
\end{itemize}

\subsection{Canli Entegrasyon Testi (Extension -> Docker Backend)}

9 Nisan 2026 sonrasindaki dogrulama turunda asagidaki adimlar canli olarak calistirildi:

\begin{enumerate}[leftmargin=1.5em]
	\item \texttt{vscode-extension} icinde TypeScript derleme: \texttt{npm run build}.
	\item Docker backend saglik kontrolu: \texttt{GET /health} \texttt{200 OK}.
	\item JWT login: \texttt{POST /api/auth/login} ile token alma.
	\item Extension key yasam dongusu simulasyonu: \texttt{POST /api/keys} ile \texttt{ca-...} anahtar uretimi.
	\item Extension cagrisi simulasyonu: \texttt{POST /v1/ask} + \texttt{X-API-Key} ile yanit alma.
	\item JWT korumali endpoint kontrolu: \texttt{GET /api/gamification/profile} + \texttt{Authorization: Bearer}.
\end{enumerate}

\subsection{Test Sonuclari}

\begin{table}[h]
\centering
\begin{tabularx}{\textwidth}{@{}lccX@{}}
\toprule
Kontrol & Beklenen & Sonuc & Not \\
\midrule
\texttt{/health} & 200 & \checkmark PASS & Backend ayakta \\
\texttt{/api/auth/login} & 200 + JWT & \checkmark PASS & Bearer token uretildi \\
\texttt{/api/keys} (create) & 201/200 + \texttt{ca-...} & \checkmark PASS & Extension key olustu \\
\texttt{/v1/ask} + \texttt{X-API-Key} & 200 & \checkmark PASS & Extension kontrati calisti \\
\texttt{/api/gamification/profile} & 200 + numeric JSON & \checkmark PASS & \texttt{progress\_percent=1.1} \\
UTF-8/JWT decode stabilitesi & 500 yok & \checkmark PASS & Encoding kaynakli crash gozlenmedi \\
\bottomrule
\end{tabularx}
\caption{Hafta 9.5 full entegrasyon test ozeti}
\end{table}

\subsection{Operasyonel Not}

PowerShell tarafinda \texttt{docker logs} ciktisinin \texttt{stderr} siniflandirmasi nedeniyle sahte \texttt{NativeCommandError} goruntulense de uygulama loglari incelendiginde bu durumun runtime hatasi degil, stream yonlendirme davranisi oldugu dogrulanmistir. Uygulama davranisi ve endpoint cevaplari test boyunca stabil kalmistir.

\section{Sonuc}

Hafta 9 ciktilariyla birlikte:
\begin{itemize}[leftmargin=1.5em]
    \item VS Code extension tarafinda calisir MVP akisina ek olarak SecretStorage ve SSE canli cikti katmani eklendi.
    \item Token satin alma mekanizmasi idempotency sertlestirmesi ile daha guvenli hale getirildi.
    \item Checkout sonrasi otomatik bakiye yenileme ile kullanici deneyimi hizlandi.
    \item Hafta 9.4 kapsaminda Agent Modu akilli yonlendirme (routing prompt + heuristik guardrail), kayip sohbet sorunu cozumu ve veritabani migrasyonu ile kullanici kontrolune dayali temiz bir mimari benimsendi.
	\item Hafta 9.5 kapsaminda extension ile Docker backend arasinda full entegrasyon testi basariyla tamamlandi; JWT/API key ayrimi, UTF-8 sanitization ve locale-safe sayisal cikti davranisi canli ortamda dogrulandi.
    \item Sistem hem urunlestirme (extension) hem gelir modeli (token ekonomisi) hem de AI kalitesi (agent routing) acisindan daha olgun ve dayanikli bir seviyeye tasindi.
\end{itemize}

\end{document}
